\begin{annexitem}{annex:p08.src025501.iso}{thm:p08-bvb501-isomorphism}{Isomorphism invariance in the source language}
\begin{proof}
Induct on the formula. Atoms use preservation of \(V\), while Boolean
constants have the same fixed truth value. Negation reverses truth on both
sides; each binary Boolean connective applies the same truth function to
the two matching induction values. For a diamond witness \(v\), relation
preservation makes \(f(v)\) a successor of \(f(w)\), and induction
preserves the argument's truth. Conversely, a target witness has a preimage
by surjectivity and relation reflection makes it a source witness.
For a box, every target successor has such a preimage and induction preserves
truth of its argument; the reverse direction follows by the same argument
for \(f^{-1}\). Thus all modality names and all constructors are covered.
\end{proof}
\end{annexitem}

\begin{annexitem}{annex:p08.src025501.game}{thm:p08-bvb501-evaluation}{Evaluation-game adequacy with all terminal and move cases}
\begin{proof}
Induct on formula complexity; each move lowers it, so infinite play is
impossible even on an infinite model. At atoms and constants the game's
declared winner is exactly the semantic one. If a disjunction is true,
Verifier chooses a true disjunct and uses its inductive winning strategy;
if it is false, each disjunct is false, so whatever she chooses Falsifier
uses the corresponding strategy. At a conjunction the roles are dual:
if both conjuncts are true Verifier answers either choice, and otherwise
Falsifier chooses a false conjunct. At negation interchange truth and
falsity and the two players, obtaining exactly the inductive result.

If a diamond is true, Verifier selects a successor witnessing its argument
and follows the winning argument strategy. If it is false, either there
is no successor and she loses immediately, or every successor's argument
is false and Falsifier has its inductive response. If a box is true,
either Falsifier has no legal move and loses, or every possible successor
has a true argument and Verifier answers with its winning strategy.
If the box is false Falsifier selects a successor with false argument and
continues its strategy. These cases prove both existence and the identity
of the winning player. For branching over an arbitrary set of successors,
choose the corresponding inductive strategies in the standing ZFC metatheory.
Expanding implication and equivalence by their truth-functional definitions
completes the redundant source syntax without adding any semantic condition.
\end{proof}
\end{annexitem}

\begin{annexitem}{annex:p08.src025501.st}{thm:p08-bvb501-standard-translation}{Truth of the standard translation}
\begin{proof}
Induct on the formula for arbitrary assignments that give \(x\) the
displayed state. At a proposition the first-order predicate has exactly
the valuation's extension. Constants have the same fixed truth values,
and every Boolean step commutes with the same classical truth function.
For a diamond the translated existential is true exactly when some state
\(v\) satisfies \(R_m(w,v)\) and the translated argument is true with
\(y\) assigned \(v\). By induction this is exactly a modal witness.
For a box the translated universal requires argument truth at precisely
all the \(R_m\)-successors of \(w\), again matching the modal clause
by induction. The freshness convention prevents capture: previously
bound variables remain distinct from each new state variable. A separate
syntactic induction shows that the only possible free variable is \(x\).
The Boolean-constant case explains why its free-variable set may be empty.
All constructors and both directions are thereby proved.
\end{proof}
\end{annexitem}

\begin{annexitem}{annex:p08.src025501.meta}{thm:p08-bvb501-modal-metatheory}{Compactness, every allowed infinite cardinal and effective enumeration}
\begin{proof}
Let \(c\) be a fresh constant and form the first-order theory
\(Q=\{\operatorname{ST}_x(\varphi)[x:=c]:\varphi\in\Sigma\}\).
Every finite subset has a model by the translation theorem, interpreting
\(c\) as its modal witness. First-order compactness
(Theorem~\ref{thm:fol-compactness}) gives a model of \(Q\).
Its distinguished state satisfies all of \(\Sigma\) by translation.

For the cardinal claim let \(\kappa\) be infinite and at least the
signature size. Add constants \(d_\xi\), \(\xi<\kappa\), and
axioms \(d_\xi\ne d_\eta\) for \(\xi\ne\eta\) to \(Q\).
Every finite subset is satisfied in the given infinite model of \(Q\),
so compactness gives a model \(A\) with at least \(\kappa\) elements.
Start with a set \(X_0\) of size \(\kappa\) containing all the
\(d_\xi^A\) and \(c^A\). At stage \(n+1\), adjoin a witness in
\(A\) for every true existential formula with parameters in \(X_n\).
There are at most \(\kappa\) such formulas and finite parameter tuples.
Choice selects the witnesses; \(X_{n+1}\) still has size \(\kappa\).
Put \(X=\bigcup_{n<\omega}X_n\). Its induced relational structure
contains the constants and is elementary in \(A\): every existential
formula with finitely many parameters from \(X\) already had those
parameters at some stage, and its witness was added at the next.
The Tarski--Vaught test
(Theorem~\ref{thm:tarski-vaught}) applies. This structure has size
exactly \(\kappa\), satisfies \(Q\), and gives the required pointed
modal model. Countable signatures allow every infinite \(\kappa\).

For effectiveness enumerate finite first-order derivations in the effective
correspondence language. Correct derivations can be checked by their finite
syntactic rules. For each modal formula \(\varphi\), compute the closed
formula \(\forall x\operatorname{ST}_x(\varphi)\), and dovetail the
search for its derivation over all modal formulas. Soundness ensures every
formula printed this way is valid by translation; first-order completeness
(Theorem~\ref{thm:fol-completeness}) ensures every valid modal formula
is eventually printed. This is semidecision, with no claimed stopping time
for an invalid input.

Finally the large-signature example does force the stated size: select for
each \(j\) a successor satisfying \(p_j\). If the same successor served
\(i\ne j\), it would violate \(\Box\neg(p_i\land p_j)\).
Hence the chosen successors are distinct. A star with one successor for
each \(j\), making only its own \(p_j\) true there, satisfies the
premises. This is a cardinal lower-bound obstruction, independent of any
search procedure for smaller models.
\end{proof}
\end{annexitem}

\begin{annexitem}{annex:p08.src025501.bisim}{thm:p08-bvb501-bisimulation-invariance}{All constructors in bisimulation invariance}
\begin{proof}
Induct on the well-founded formula construction. Atomic harmony supplies
the atomic case and constants are fixed. Negation and every Boolean
connective preserve equality of their argument truth values. For a diamond
true at \(w\), take its witness \(v\); zig supplies a related target
witness and induction preserves argument truth. Zag supplies the reverse
direction. For a box true at \(w\), every target successor has a related
source successor by zag, whose true argument transfers by induction;
the reverse uses zig. For any set-sized conjunction or disjunction apply
induction to each operand and the semantic universal or existential clause
for the connective. This is transfinite structural induction, not an
assumption that every infinitary formula has finite depth.
\end{proof}
\end{annexitem}

\begin{annexitem}{annex:p08.src025501.transform}{thm:p08-bvb501-transformations}{The four source constructions and the restriction counterexample}
\begin{proof}
Within a tagged disjoint union the component inclusion preserves atoms.
Every successor of an included point stays in that component, so the graph
of the inclusion satisfies both transition clauses. In a generated
restriction the same argument uses successor closure; the restricted
relations and valuations supply the other direction and atomic harmony.
For a bounded morphism its graph has atomic harmony. Relation preservation
supplies zig and the stipulated back condition supplies zag. Surjectivity
is unnecessary for truth at an image point, though required for calling
the entire target a bounded morphic image. Apply the invariance theorem
to all three graphs.

For an unraveled path its successors are exactly the paths extending it
by one labelled edge from its last state. Last-state projection therefore
preserves relations and answers every original successor with that extension;
its valuation clause gives atomic harmony. Its image is precisely the
reachable generated domain, including the length-zero path. This proves
both truth preservation and the stated image. An unreachable original
state is not its image, so surjectivity onto an arbitrary original model
cannot be inferred. Finally a root with one terminal successor satisfies
\(\Diamond\topv\); restricting to the root makes that formula false.
This restriction is induced but not generated, proving the boundary.
\end{proof}
\end{annexitem}

\begin{annexitem}{annex:p08.src025501.contract}{thm:p08-bvb501-contraction}{Equivalence, well-defined quotient and absence of further collapse}
\begin{proof}
The identity is a self-bisimulation. Inverses of bisimulations are
bisimulations because the two transition clauses exchange places.
Nonempty compositions are also bisimulations: atomic agreement is transitive,
and a successor can first be matched through one relation and then through
the other; reversing the steps proves back. Any nonempty union satisfies
the clauses by selecting a constituent containing the pair under consideration.
Consequently the union \(\sim\) of all self-bisimulations is reflexive,
symmetric and transitive: for transitivity compose constituents witnessing
the two pairs. Atomic harmony makes the quotient valuation independent
of the representative.

Projection preserves an original edge by the definition of the quotient
edge. Conversely, a quotient edge out of \([w]\) has witnesses
\(a\sim w\), \(aR_mb\) and \([b]=[v]\).
The bisimulation \(\sim\) supplies \(wR_mc\) with \(c\sim b\).
Thus projection of \(c\) is the specified quotient successor. This is
the back clause, and projection is plainly onto the classes.

If a quotient bisimulation related two different classes, compose it with
projection's graph and its inverse. This is a bisimulation on the original
model relating representatives of the two classes. It lies in \(\sim\),
contradicting their difference. If a further bounded morphism identified
two quotient states, its graph composed with its inverse would similarly
give a quotient self-bisimulation relating them. Thus no such identification
is possible. This proves the claimed sense of removing bisimulation
redundancy without identifying it with arbitrary graph isomorphism.
\end{proof}
\end{annexitem}

\begin{annexitem}{annex:p08.src025501.hm}{thm:p08-bvb501-hennessy-milner}{Finite successor sets give the modal converse}
\begin{proof}
Let \(E\) relate modally equivalent states. Atomic harmony is immediate.
Suppose \(wEw'\) and \(wR_mv\). If \(w'\) has no \(m\)-successor,
\(\langle m\rangle\topv\) already contradicts equivalence. Otherwise
list its finitely many successors \(u'_1,\ldots,u'_r\).
If none is equivalent to \(v\), choose for each one a formula
\(\psi_i\) true at \(v\) and false at \(u'_i\), negating a
distinguishing formula if its orientation is reversed.
Then \(\langle m\rangle\bigwedge_{i=1}^r\psi_i\) is true at
\(w\) and false at \(w'\), a contradiction. Hence there is a matching
successor in \(E\). Symmetrically apply finiteness of the source
successor set to obtain zag. This argument works for each modality and
uses only a finite conjunction even with infinitely many proposition
letters. The nonempty relation containing the initial pair is a bisimulation.
\end{proof}
\end{annexitem}

\begin{annexitem}{annex:p08.src025501.lift1}{thm:p08-bvb501-first-lifting}{Saturated extensions, successor realization and the first lift}
\begin{proof}
For a relational structure \(A\), add names for all its elements and
take its elementary diagram. For each complete one-variable type over
each finite parameter tuple consistent with its theory, add a fresh
realizer constant and all formulas of that type. Every finite subset of
this enlarged theory is satisfiable in \(A\): each finite fragment of
a type has a realizing witness there, since its existential closure is
consistent with the complete theory of the named parameters and hence
true. Different realizer constants have no mutual constraints, so their
finite requirements can be witnessed together. Compactness supplies an
elementary extension realizing the entire family. There is a set of types
here, since the language and domain are sets.

Repeat this at each finite stage and take the elementary-chain union.
The union is elementary over every stage: induction on formulas transfers
an existential witness found at a later stage back to an earlier stage
containing its finite parameter tuple; the converse preserves witnesses
forward. Every finite tuple of union parameters belongs to some stage.
A complete one-type over it consistent with the union is consistent with
that stage's identical theory and is realized at the next stage. Hence
the union is \(\omega\)-saturated for one-types. For a finite-tuple type,
realize its first-coordinate marginal, then its next-coordinate conditional
type over that coordinate and the original parameters, and continue.
Consistency at each step follows because the already realized tuple has
the complete marginal type. Thus all finite-tuple types are realized.

In an \(\omega\)-saturated correspondence structure, a finitely
realizable successor condition \(\Gamma\) translates to the partial
first-order type
\(\{R_m(w,y)\}\cup\{\operatorname{ST}_y(\gamma):\gamma\in\Gamma\}\)
over the parameter \(w\). Compactness extends it to a complete type;
saturation realizes it. This proves modal saturation.

On two modally saturated models let \(E\) be modal equivalence.
For \(wEw'\) and \(wR_mv\), take the complete modal theory of \(v\),
including the negation of every formula false there. Every finite part has
an \(m\)-successor witness at \(w'\): its conjunction is true at
\(v\), so its diamond is true at \(w\), and therefore at \(w'\).
Modal saturation realizes the entire theory at a successor \(v'\),
giving \(vEv'\). The symmetric argument proves zag. Atomic harmony
is immediate. Thus modal equivalence is a bisimulation on these models.

If the original pointed models are modally equivalent, extend them by the
construction above. Elementarity and the standard translation preserve
their full modal theories, so their named points are still equivalent and
the preceding argument gives a bisimulation. Conversely, bisimilar
elementary extensions agree modally by invariance, and elementarity transfers
that agreement to the original points. Both directions of the first lift
are proved; it is not a substitution for the source's second lift.
\end{proof}
\end{annexitem}

\begin{annexitem}{annex:p08.src025501.character}{thm:p08-bvb501-characterization}{The all-model characterization with compactness and saturation}
\begin{proof}
A standard translation is bisimulation-invariant by translation truth
and modal invariance. For the converse let \(C\) consist of all modal
consequences of \(\alpha(x)\), written as their standard translations.
We first prove \(C\entails\alpha\). Suppose \((M,w)\) satisfies
\(C\). The first-order theory asserting \(\alpha(c)\), together
with translations at \(c\) of the complete modal theory of \((M,w)\),
is finitely satisfiable. Otherwise some finite conjunction \(\theta\)
of formulas true at \(w\) is inconsistent with \(\alpha\).
Then \(\neg\theta\) is a modal consequence of \(\alpha\) in
\(C\), contradicting truth of \(\theta\) at \(w\).
Compactness gives \((N,v)\) satisfying \(\alpha\) and the same
complete modal theory as \((M,w)\).

By the first lifting theorem, there are bisimilar pointed elementary
extensions \((M^+,w)\), \((N^+,v)\). The latter satisfies
\(\alpha\) by elementarity. Bisimulation invariance transfers it to
the former, and elementarity transfers it back to \((M,w)\).
This proves \(C\entails\alpha\). Compactness of
\(C\cup\{\neg\alpha(c)\}\) selects a finite subfamily whose
conjunction implies \(\alpha\). By definition \(\alpha\) implies
each member, so it is equivalent to that conjunction of modal translations,
itself the translation of one modal formula. Empty conjunctions are
\(\topv\); an inconsistent \(\alpha\) is covered by \(\bottom\).
The proof used only set-sized theories and is valid for arbitrary signatures
in the standing ZFC metatheory. Compactness over finite models is not
available, so no finite-model characterization is claimed from this proof.
\end{proof}
\end{annexitem}

\begin{annexitem}{annex:p08.src025501.game-infinite}{thm:p08-bvb501-infinite-game}{Strategies, largest bisimulation and ordinal descent}
\begin{proof}
Given a bisimulation, Duplicator keeps the current pair in it by choosing
the zig or zag witness. Atomic disagreement cannot occur and every required
response exists. If no Spoiler move exists she wins immediately; otherwise
she survives an infinite play. Conversely, from a winning strategy take
all state pairs reached by plays following it. They agree atomically,
and after each possible Spoiler move the strategy supplies a related
response. Thus this nonempty relation satisfies both bisimulation clauses.
The correspondence is a conversion of witnesses, not a claim that every
history-dependent strategy corresponds to one uniquely determined relation.

To obtain a Spoiler strategy when no bisimulation exists, define \(Z_0\)
to be all cross-model pairs. Let \(Z_{\beta+1}\) consist of pairs
agreeing on atoms whose every transition on either side has a matching
pair in \(Z_\beta\); at limits intersect the earlier relations.
The sequence decreases and stabilizes, since strict stages remove distinct
pairs from a set. Its fixed point is a bisimulation if nonempty, and every
bisimulation is contained in every stage by transfinite induction. Hence
it is the greatest relation satisfying the clauses.

For a pair outside that fixed point, its first exclusion occurs at a
successor stage \(\beta+1\): a limit cannot be its first exclusion
because that stage is an intersection. Atomic disagreement gives an
immediate win. Otherwise failure of a transition clause supplies a Spoiler
move for which every possible response is outside \(Z_\beta\).
An impossible response wins immediately. For a response that exists, its
first-exclusion ordinal is strictly smaller than \(\beta+1\).
Repeat. There is no infinite strictly descending sequence of ordinals,
so play terminates in a Spoiler win. This proves determinacy without
assuming a uniform finite bound on the lengths of the plays.
\end{proof}
\end{annexitem}

\begin{annexitem}{annex:p08.src025501.game-finite}{thm:p08-bvb501-finite-game}{Finite signatures, finite branching and the actual failure case}
\begin{proof}
First a distinguishing formula of depth at most \(k\) always supplies
a Spoiler strategy. For a Boolean difference choose a differing operand,
and if necessary negate it to orient truth. At an atomic difference win
immediately. For a diamond true on one side but false on the other,
choose its witnessing successor. Every possible opposing successor falsifies
the argument, so continue with that difference at one lower depth. A box
is the dual diamond case. This process uses at most the modal depth in
rounds, whatever the branching or signature size.

For a finite relevant signature build finite-depth descriptions. A depth-zero
description is the conjunction of the true atoms and negations of the false
atoms. There are finitely many such patterns. Recursively a depth-\(j+1\)
description records that atomic pattern and, for each modality, the set
\(S_m\) of depth-\(j\) descriptions occurring at its successors.
There are finitely many records. Its characteristic formula is the atomic
description conjoined, for every \(m\), with
\[
 \bigwedge_{\rho\in S_m}\langle m\rangle\chi_\rho
 \quad\land\quad
 [m]\bigvee_{\rho\in S_m}\chi_\rho.
\]
The empty conjunction and disjunction are \(\topv\) and \(\bottom\).
Induction proves that this formula describes exactly the record. Two equal
depth-\(j+1\) records permit matching a chosen successor with the same
depth-\(j\) record. Different records either already differ on an atom
or have a modality with an unmatched successor record. Induction therefore
identifies equal records with Duplicator wins, and different records with
Spoiler wins and a characteristic distinguishing formula. This proves
the first equivalence for arbitrary models over that finite signature.

For arbitrary signatures on image-finite models argue by induction on \(k\).
A zero-round loss has an atomic witness. At a positive-round loss with
atomic agreement, Spoiler has a successor move such that every response
loses in at most \(k-1\) remaining rounds. There are finitely many
responses for that chosen modality. Induction selects one oriented formula
true at the chosen successor and false at each particular response. Their
finite conjunction, preceded by the chosen diamond, distinguishes the
initial points and has depth at most \(k\). If there are no responses,
use \(\langle m\rangle\topv\). The already proved other direction
completes this equivalence. Induction on the finite game also proves its
determinacy by inspecting each legal Spoiler move and its possible replies.

In the one-round counterexample, a finite formula's leaf truth depends
only on its finite atom set, because every leaf has no successors.
Every pattern on that set occurs among the finite-subset leaves, so adding
the all-atoms leaf changes neither existence nor universality of argument
truth at successors. Structural induction makes all root formulas agree.
The first-round move to the new leaf has no atomically agreeing finite-subset
response, proving the asserted failure of the converse without the
conditions. It also proves these countable models are not bisimilar.

For the Figure~14 chain example, for any fixed \(k\) Duplicator can
match a move into a finite chain with its exact counterpart. A move into
the infinite chain is answered by a finite chain with at least \(k\)
steps available. All subsequent moves advance one step, so she survives
the remaining rounds. This holds relative to every finite atom signature,
with every atom false, and proves finitary modal equivalence. In the
unlimited game Spoiler chooses the infinite chain and advances on it;
every finite response eventually ends and cannot answer a further move.
Thus there is no bisimulation. The two counterexamples have different
roles: one refutes uniform finite-round formula adequacy with infinitely
many atoms, and the other refutes unrestricted Hennessy--Milner equivalence
and shows why winning unlimited play need not have a fixed finite bound.
\end{proof}
\end{annexitem}

\begin{annexitem}{annex:p08.src025501.modal-infinitary}{thm:p08-bvb501-infinitary-modal}{A characteristic infinitary modal formula}
\begin{proof}
Start a self-model approximation \(A_0\) with the atomically agreeing
pairs; let \(A_{\gamma+1}\) contain those pairs in \(A_0\) whose
transitions match through \(A_\gamma\), and intersect at limits.
The ordinal-descent method in the game proof gives a stabilization stage
\(\alpha\). For every \(a\in W\)
define formulas \(\chi_a^\gamma\) by transfinite recursion. At zero take
the complete conjunction of its true atoms and the negations of its false
atoms. At a successor take that conjunction, together with, for every \(m\),
\[
 \bigwedge_{aR_mb}\langle m\rangle\chi_b^\gamma
 \quad\land\quad [m]\bigvee_{aR_mb}\chi_b^\gamma.
\]
At a limit conjoin all preceding formulas for \(a\). Empty successor
sets use the Boolean constants as in the finite-depth construction.
Induction shows that \(N,t\entails\chi_a^\gamma\) exactly when
\((a,t)\) belongs to the corresponding cross-model \(A_\gamma\).
In particular the self-model stabilization gives
\(M\entails\chi_a^\alpha\to\chi_a^{\alpha+1}\) for every \(a\).

Let \(H=\bigwedge_{a\in W}(\chi_a^\alpha\to\chi_a^{\alpha+1})\).
Define \(\delta_{M,w}\) as \(\chi_w^\alpha\) conjoined with
\([m_1]\cdots[m_j]H\) for every finite modality string, including
the empty string. The formula is true at \(M,w\) by stabilization.
If \(N,t\) satisfies it, put \(aEb\) when \(b\) is reachable
from \(t\) and satisfies \(\chi_a^\alpha\). Its root pair is present.
At every reachable \(b\), the \(H\) conjuncts imply
\(\chi_a^{\alpha+1}\). The successor definition then supplies both
transition matches in \(E\), and atomic harmony comes from the
zero-stage conjunct. Matched target successors stay reachable. Thus \(E\)
is a bisimulation. Conversely, bisimulation invariance for infinitary
formulas transfers \(\delta_{M,w}\) from \(M,w\) to any bisimilar
point. This proves its exact characteristic property.

All conjunction indices are sets. If the model and signature are countable,
the self-pair set is countable and the descending relation stabilizes at
a countable ordinal: before stabilization each strict stage removes a
distinct pair, so there cannot be \(\omega_1\) strict stages.
All recursion stages and all conjunctions in this construction are then
countable. Agreement on countable-infinitary modal formulas implies truth
of the characteristic formula and hence bisimilarity; the reverse is
invariance. Countable connective indices still do not force finitary syntax.
\end{proof}
\end{annexitem}

\begin{annexitem}{annex:p08.src025501.potential}{thm:p08-bvb501-potential-isomorphism}{Infinitary invariance and a characteristic sentence}
\begin{proof}
For invariance use structural induction on an infinitary formula at any
tuple matched by \(f\in I\). Atomic preservation is the hypothesis.
Boolean and set-sized connective steps apply the induction hypothesis to
every operand. For an existential witness in \(M\), forth extends \(f\)
to that witness and induction gives the matching truth in \(N\). Back
gives the converse. Universal quantifiers follow by negation. Finite free
variable sets ensure the tuple remains contained in a finite map throughout
this induction. No fixed finite formula rank is assumed.

For the characteristic sentence construct descending relations \(B_\gamma\)
on all pairs of same-length finite tuples of \(M\). At zero require
agreement on all atomic formulas. At a successor retain the preceding
relation and require that each one-element extension of either tuple has
a matching extension in the preceding relation. At limits intersect.
This process stabilizes at some ordinal \(\alpha\), because there is
a set of finite tuple pairs and strict stages remove distinct pairs.

For each finite tuple \(\bar a\) define \(\theta_{\bar a}^\gamma\).
At zero take its complete positive and negative atomic diagram. Put
\[
\begin{split}
\theta_{\bar a}^{\gamma+1}(\bar x)={}&
 \theta_{\bar a}^\gamma(\bar x)
 \land\bigwedge_{b\in M}\exists y\,
       \theta_{\bar a b}^\gamma(\bar x,y)\\
 &\land\forall y\,\bigvee_{b\in M}
       \theta_{\bar a b}^\gamma(\bar x,y),
\end{split}
\]
and conjoin the earlier stages at limits. Induction shows that these
formulas express the corresponding tuple relation against any other
structure. Repeated tuple entries are handled by the atomic equality
diagram. Define
\[
\delta_M=\theta_{()}^\alpha\land
 \bigwedge_{n<\omega}\ \bigwedge_{\bar a\in M^n}
 \forall\bar x\,
 (\theta_{\bar a}^\alpha(\bar x)
       \to\theta_{\bar a}^{\alpha+1}(\bar x)).
\]
Stabilization of the self-model relation makes this sentence true in \(M\).
In any model \(N\) of it, collect the finite tuple maps whose target
tuples satisfy \(\theta_{\bar a}^\alpha\). Atomic diagrams make them
well-defined partial isomorphisms. The empty tuple supplies a map, and
the successor clauses, forced by \(\delta_M\), supply both extension
conditions. Thus they form a potential isomorphism. Conversely invariance
transfers \(\delta_M\) to every potentially isomorphic structure.
Equality of full infinitary theories therefore supplies truth of this
sentence and gives precisely the converse claimed.

If \(M\) and the signature are countable, the tuple-pair set is countable,
the first stabilization ordinal is countable, and every atomic diagram
and connective index in this construction is countable. Thus the sentence
belongs to \(L_{\omega_1\omega}\). For countable \(N\), start with
a map in the family and alternately include the next element of each
enumerated domain and range. The union is a bijection preserving all
atomic formulas and hence an isomorphism. This proves the countable
characterization without asserting that an infinitary sentence forces
countability of its models.

The finite-structure comparison in the body also follows directly.
Elementary equivalence fixes the finite domain size. If no bijection were
an isomorphism, choose for each of the finitely many bijections one atomic
fact it violates. An existential sentence listing the whole finite domain
and conjoining that finite diagram is true on one side and false on the
other, a contradiction. The rational/real order example satisfies the
extension conditions by picking points in the appropriate nonempty intervals
or rays; its different cardinalities rule out an isomorphism.
\end{proof}
\end{annexitem}

\begin{annexitem}{annex:p08.src025501.assignments}{thm:p08-bvb501-assignment-bridge}{The assignment-reset correspondence with its existential scope}
\begin{proof}
Suppose \(I\) is a restriction-closed potential isomorphism. Pick a map
containing some \(a_0\mapsto b_0\), possible by forth and nonempty
domains. Consider assignments \(g,h\) finitely different from the
constant assignments with values \(a_0,b_0\), respectively. Relate them
when the finite coordinate graph
\(\{(g(x),h(x)):x\text{ a variable}\}\) is a map in \(I\).
The constant assignments are related. Every atomic proposition agrees
because that map preserves all atomic first-order formulas on their values.

If a move resets \(g(x)\) to \(a\), extend the coordinate map through
forth to include \(a\mapsto b\). Answer by resetting \(h(x)\) to
\(b\). The new coordinate graph is a restriction of this extension,
so belongs to \(I\), and both assignments still have finite support.
Back supplies the reverse move condition. Hence this is a nonempty
bisimulation. In particular its existence does not require a relation
covering the assignment models' entire, possibly enormous, domains.

Conversely take a nonempty assignment-model bisimulation. Collect every
finite coordinate restriction of every related assignment pair. Atomic
equality agreement makes each such restriction a well-defined injective
map; agreement of all atomic proposition letters makes it a partial
isomorphism. Given one restriction and any \(a\in M\), choose a
variable outside its finite coordinate list and reset it to \(a\).
The bisimulation supplies a reset on the other side while all listed
coordinates remain unchanged. The new restriction extends the old map
to include \(a\). A reverse reset includes any prescribed range element.
This proves both extension conditions and produces the potential
isomorphism. Atoms with constants or compound terms cause no omission:
they are themselves among the assignment model's proposition letters.
\end{proof}
\end{annexitem}

\begin{annexitem}{annex:p08.src025501.fo-game}{thm:p08-bvb501-first-order-game}{The finite relational signature rank theorem}
\begin{proof}
A winning \(k\)-round tuple strategy preserves all formulas of quantifier
rank at most \(k\) by induction: atoms use the initial partial
isomorphism, Boolean constructors preserve agreement, and an existential
witness is played as the next game move. The response preserves its
matrix by induction with one fewer round. The opposite direction uses
the same argument for moves in the other structure; universal formulas
are dual.

For the converse, finite relational vocabulary gives finitely many
atomic diagrams on a fixed finite tuple of variables, including the
finitely many constants. There are no unbounded function terms here.
Define a rank-zero record by this atomic diagram. Recursively a rank-
\(j+1\) record contains the atomic diagram and the set of rank-
\(j\) records of tuples obtained by appending one arbitrary element.
For the finitely many tuple lengths needed in a \(k\)-round game,
induction gives finitely many records. A characteristic formula consists
of its atomic diagram and, for that extension-record set \(S\),
\[
 \bigwedge_{\rho\in S}\exists y\,\chi_\rho(\bar x,y)
 \quad\land\quad
 \forall y\,\bigvee_{\rho\in S}\chi_\rho(\bar x,y).
\]
Its quantifier rank is at most \(j+1\). Induction verifies its exact
record description. Equal records let Duplicator answer each move with
an element giving the same lower-rank record; different records give
a characteristic formula distinguishing the tuples. Thus agreement on
all formulas of rank at most \(k\) forces equal records and the desired
strategy. This proof requires finitary descriptions of the zero and
successor records; it supplies no unrestricted strategy-to-finitary-formula
converse for an infinite atomic vocabulary.
\end{proof}
\end{annexitem}

\begin{annexitem}{annex:p08.src025501.interpolate}{thm:p08-bvb501-relative-interpolation}{Common-signature interpolation by compactness}
\begin{proof}
A modal interpolant gives entailment along common-signature bisimulations
by modal invariance. Conversely, let \(C\) contain all basic modal
\(L\)-formulas whose standard translations follow from \(\alpha\).
Suppose a pointed full \(\beta\)-signature model \(M,w\)
satisfies \(C\) and \(\neg\beta\). Add a constant \(c\) and
combine \(\alpha(c)\) with the translated complete modal
\(L\)-theory of \(M,w\). Every finite part is consistent: if
\(\alpha\) conflicted with a finite conjunction \(\gamma\)
of that modal theory, \(\neg\gamma\) would belong to \(C\),
contrary to truth of both \(C\) and \(\gamma\) at \(w\).
Compactness therefore gives a full \(\alpha\)-signature model
\(N,v\) satisfying \(\alpha\) and modally agreeing with
\(M,w\) on \(L\).

Take saturated elementary extensions in their respective full signatures,
by Theorem~\ref{thm:p08-bvb501-first-lifting}. Their \(L\)-reducts
are modally saturated: a successor type in the reduct is also a type
in the full signature with its finite parameters and accessibility
condition retained. Common modal equivalence supplies an \(L\)-
bisimulation between the named points. Elementarity retains
\(\alpha\) on one side and \(\neg\beta\) on the other,
contradicting the assumed structural entailment. Hence
\(C\entails\beta\). Compactness, applied to its standard
translations with \(\neg\beta(c)\), supplies a finite
\(C_0\subseteq C\) whose conjunction \(\theta\) implies
\(\beta\). Every member of \(C_0\) follows from
\(\alpha\), so \(\theta\) is the desired interpolant.
An empty conjunction is \(\topv\). This proof uses all-model
compactness, not finite-model compactness, and does not check private
symbols in the bisimulation.
\end{proof}
\end{annexitem}

\begin{annexitem}{annex:p08.src025501.modal-interpolate}{thm:p08-bvb501-modal-interpolation}{Private valuations on related pairs}
\begin{proof}
Let \(M,w\) and \(N,v\) have a bisimulation \(E\) checking
all fixed modalities and just the common proposition letters. Suppose
\(M,w\entails\varphi\) and \(N,v\not\entails\psi\).
Make a model with domain \(E\subseteq W_M\times W_N\), and
put \((a,b)R_m(a',b')\) iff \(aR_m^Ma'\),
\(bR_m^Nb'\), and \(a'Eb'\). Give each atom of
\(\varphi\) its first-coordinate valuation and each atom of
\(\psi\) its second-coordinate valuation. Common atoms agree
by harmony; other atoms can be arbitrary. Each projection preserves
edges, and its back condition follows from the appropriate forth or
back clause of \(E\). On the relevant atom signature it preserves
valuations too. Bounded-morphism truth preservation gives both
\(\varphi\) and \(\neg\psi\) at \((w,v)\), contrary
to \(\varphi\entails\psi\). Thus the standard translations
satisfy entailment along these common-atom bisimulations. The preceding
theorem gives a common-atom modal interpolant; translation correctness
gives its two modal entailments. Both projections check all fixed
modalities; this argument does not treat private relation symbols as
private valuations.
\end{proof}
\end{annexitem}

\begin{annexitem}{annex:p08.src025501.check}{thm:p08-bvb501-model-checking}{Truth sets and graph scans}
\begin{proof}
Induct in subformula order. Valuations and Boolean constants give the
right base sets; Boolean set operations agree with their truth clauses.
For diamond, scan backwards from the argument's computed set: exactly
states with a successor in that set are marked. For box, begin with all
states and remove any state having an edge to a point outside the
argument set. This is precisely universal successor truth, including
vacuous truth at dead ends. Thus every stored truth set is correct.
A Boolean step scans at most \(|W|\) bits and a modal step at most
\(|W|+|R_m|\) adjacency entries. There are at most
\(s(\varphi)\) steps; storing every answer set uses
\(O(s(\varphi)|W|)\) space. A matrix for one relation has
\(|W|^2\) entries, giving the source's coarser bound. These
counts assume the explicit graph, not a succinct description of its edges.
\end{proof}
\end{annexitem}

\begin{annexitem}{annex:p08.src025501.until}{thm:p08-bvb501-until-check}{Least predecessor closure is path Until}
\begin{proof}
Write \(A=[\![\psi]\!]\), \(B=[\![\theta]\!]\).
The queue computes the least set \(X\) containing \(A\) and
closed under adding a \(B\)-state with a successor in \(X\).
Each initial mark has a length-zero witness path. Each later mark
prepends its edge to an already witnessed path, proving soundness.
Conversely induction on a witness path's length shows its endpoint is
initially marked and processing its marked suffix eventually marks
each predecessor. Thus all and only path-Until states are marked.
Each state enters the queue once and its incoming list is scanned once,
giving \(O(|W|+|R|)\) time. In a directed acyclic unique-path
graph, the points with \(sR^*u\) and \(uR^+t\) are exactly
the unique path's points before its endpoint: concatenated paths through
\(u\) must be that unique path, and acyclicity excludes the endpoint.
Hence the interval and path clauses agree there. The body's merging-
branch example proves the restriction necessary for the algorithm's
claimed interval interpretation on arbitrary irreflexive graphs.
\end{proof}
\end{annexitem}

\begin{annexitem}{annex:p08.src025501.two-variable}{thm:p08-bvb501-two-variable}{Scoping in the alternating translation}
\begin{proof}
Induct simultaneously that the \(x\)-translation has at most free
variable \(x\), uses only \(x,y\), and is true at an assigned
\(x\)-value \(w\) iff \(M,w\entails\sigma\); prove the
corresponding statement with \(x,y\) reversed. Atomic and Boolean
cases are immediate. In an \(x\)-diamond clause the outer binding
of \(y\) selects an \(R_m\)-successor of the assigned \(x\)
value, where the \(y\)-inductive hypothesis applies. Any internal
binding of \(x\) affects only that translated argument, not the
outer edge atom. The universal box clause uses the same scoping and
the other two clauses are symmetric. This proves both truth preservation
and the free-variable property. Existential closure at the initial point
therefore preserves satisfiability in both directions.
\end{proof}
\end{annexitem}

\begin{annexitem}{annex:p08.src025501.filtration}{thm:p08-bvb501-filtration}{The paired minimum--maximum filtration}
\begin{proof}
Each displayed constraint depends only on the two classes, since its
modal formula and argument belong to \(\Sigma\). If an original
edge \(aR_mb\) connects the classes of \(w,v\), truth of an
argument at \(v\) transfers to \(b\), producing diamond truth
at \(a\), hence at \(w\). Truth of a box at \(w\) transfers
to \(a\), its argument to \(b\), then to \(v\). Thus every
minimum edge satisfies both maximum constraints, proving
\(R_m^s\subseteq R_m^l\).

Induct on \(\sigma\in\Sigma\). Atoms use inherited valuations;
Boolean cases use subformula closure. An original diamond witness edge
survives in every intervening relation, and induction preserves its
argument truth. Conversely a quotient diamond witness has an argument
true in the original model by induction; the diamond constraint gives
the original diamond. An original true box forces its argument at
every quotient successor by the box constraint and induction. An
original false box has an original counterexample edge; that edge
survives and induction retains the false argument in the quotient.
These prove both modal cases. Every class has a distinct
\(\Sigma\)-truth pattern, giving at most \(2^{|\Sigma|}\)
classes. No duality expansion has enlarged the counted subformula set.
\end{proof}
\end{annexitem}

\begin{annexitem}{annex:p08.src025501.finite-decision}{thm:p08-bvb501-finite-decision}{Filtration gives a terminating exhaustive search}
\begin{proof}
Filter a satisfying pointed model through the input's subformulas.
There are \(s(\varphi)\) of these; the minimum paired filtration
keeps the input true in at most \(2^{s(\varphi)}\) states.
For each domain size from one to that bound there are finitely many
binary relations for each relevant modality and valuations for each
relevant atom. Enumerate all these models and apply the truth-set
algorithm. A successful check proves satisfiability; exhausting the
finite search proves unsatisfiability by the bound. Nonoccurring symbols
can be interpreted arbitrarily, including in an infinite ambient
signature. Validity is unsatisfiability of the negated input, decided
with that input's own subformula bound.
\end{proof}
\end{annexitem}

\begin{annexitem}{annex:p08.src025501.tree-depth}{thm:p08-bvb501-tree-depth}{Truncating a finitely branching unraveling}
\begin{proof}
In the rooted tree unraveling, truncate after level \(d\).
Induct on formulas that at a retained node of level \(j\), every
formula of depth at most \(d-j\) keeps its truth. Boolean cases
are immediate. In a modal case \(j<d\), every original immediate
successor is retained, and its argument has depth at most
\(d-j-1\), so induction applies to all successors. At level
\(d\) only depth-zero formulas are covered; no modal truth is
asserted there. At the root this is the required preservation.
For a satisfiable input first take its finite filtration model and
unravel from its satisfying point. Bounded-morphism preservation keeps
root truth. With the finite relevant modality signature the tree has
finite branching; its depth-\(\operatorname{md}(\varphi)\)
truncation is finite and still satisfies the input. Neither the untruncated
unraveling nor arbitrary deletion in an unexpanded graph would prove
this finite-tree statement.
\end{proof}
\end{annexitem}

\begin{annexitem}{annex:p08.src025501.space-search}{thm:p08-bvb501-space-search}{Depth-first truth assignments with polynomial storage}
\begin{proof}
Let \(\Sigma\) be the input's subformula set. At remaining depth
\(d\), choose one truth bit for each member of modal depth at most
\(d\), respecting its Boolean truth table and all obligations
received from the parent. At the root require the input bit true.
Primitive implication and equivalence use their truth tables directly,
avoiding an exponentially duplicating expansion into formula trees.

For each modality, all boxes marked true impose their arguments at
every successor; all diamonds marked false impose their arguments
false at every successor. Each diamond marked true demands a successor
with its argument true, and each box marked false demands a successor
with its argument false. Recursively check one successor per demand
with those universal obligations and that witness obligation, at depth
\(d-1\). Different demands may use different successors. Use no
successors when there is no demand. At depth zero only Boolean and
atomic consistency remain. Reject inconsistent choices or failed children.

Induction on \(d\) proves that accepting choices construct a finite
tree satisfying exactly their local bits: every existential requirement
has its witness and every universal requirement holds at all constructed
successors. Conversely a satisfying depth-bounded tree supplies actual
local truth bits and witnesses making all recursive checks accept.
The preceding tree theorem supplies such a tree whenever the input
is satisfiable. Thus the procedure decides both satisfiable and
unsatisfiable inputs. A frame stores \(O(|\Sigma|)\) bits and the
index of its next demand. The stack has at most modal depth plus one
frames; discard a processed child before checking the next. Storage is
therefore polynomial, though the number of branches can be exponential.
The deterministic conversion uses the general NPSPACE--PSPACE equality,
Theorem~\ref{thm:p04-savitch}, not a modal hardness argument.
\end{proof}
\end{annexitem}

\begin{annexitem}{annex:p08.src025501.finite-bisim-algorithm}{thm:p08-bvb501-finite-bisimulation-algorithm}{Deletion computes the greatest bisimulation}
\begin{proof}
Every bisimulation begins within the atomically agreeing pairs. If all
its pairs survive a stage, its matching successor witnesses survive
too, so none of its pairs can be deleted at the next stage. Induction
therefore places every bisimulation within every stage. At stabilization,
each surviving pair agrees atomically and has its required forth and
back witnesses in the surviving relation. Hence the stabilized relation
is itself a bisimulation when nonempty and contains every bisimulation.
A designated pair survives iff a bisimulation contains it; an empty
table correctly rejects all pairs. There are \(|W_M||W_N|\)
initial pairs and every nonstationary round deletes at least one.
Scanning pairs, modalities, successors and possible responses against
the current finite table has polynomial cost in the explicit input;
the bounded number of scans remains polynomial. This needs neither
an optimal refinement algorithm nor formula enumeration.
\end{proof}
\end{annexitem}

\begin{annexitem}{annex:p08.src025501.e-collapse}{thm:p08-bvb501-printed-e-collapse}{The duplicated E argument collapses the extension to Id}
\begin{proof}
Let E' be the schema with its final argument \(\varphi\) instead
of \(\psi\). T and its negated substitutions give both
\(p\to\Diamond p\) and \(\Box p\to p\). Substituting
\(\topv\) for \(\psi\) in E' gives
\(\Diamond\Diamond p\to\Box(\Diamond p\lor\Box p)\).
T applied to \(\Diamond p\) gives
\(\Diamond p\to\Diamond\Diamond p\). Also
\(\Box p\to p\to\Diamond p\), so by necessitation and K,
\(\Box(\Diamond p\lor\Box p)\to\Box\Diamond p\).
Together these derive \(\Diamond p\to\Box\Diamond p\).
Substitute \(\neg p\) and contrapose using duality to obtain
\(\Diamond\Box p\to\Box p\). M and the first of these
implications now give
\(\Diamond p\to\Box\Diamond p\to\Diamond\Box p\to\Box p\).
T therefore gives \(p\to\Box p\), and its other form gives
\(p\leftrightarrow\Box p\). Uniform substitution supplies every
Id axiom, so Id is contained in this extension.

Conversely every T, M, E' and Q instance is valid on identity frames.
In Id, induction on formulas proves equivalence to their Boolean
erasures, replacing every box and diamond by its argument. A formula
valid on identity frames has a classically valid erasure; propositional
completeness and those equivalences prove it in Id. Thus all the
extension's axioms and theorems are in Id. The two logics are equal,
and this same erasure argument proves their identity-frame completeness.
Hence the duplicated argument cannot support the cited incompleteness
claim. Only theorem necessitation, not necessitation of a premise,
is used in this calculation.
\end{proof}
\end{annexitem}

\begin{annexitem}{annex:p08.src025501.correspondents}{thm:p08-bvb501-elementary-correspondents}{Valuations witnessing each necessary condition}
\begin{proof}
For transitive \(R\), a box-true point has its argument true at every
successor of a successor, hence satisfies the iterated box. If transitivity
fails with \(wRu,uRv,\neg wRv\), set \(V(p)=R[w]\).
Then \(\Box p\) is true at \(w\) while \(\Box\Box p\)
fails through \(u,v\). For a partial function a diamond witness
is the only possible successor, hence box truth follows. Conversely
two distinct successors \(u,v\) are separated by \(V(p)=\{u\}\),
making the diamond true and the box false at their source.

Identity frames validate \(p\leftrightarrow\Box p\).
If \(w\) has a distinct successor \(v\), the valuation
\(V(p)=\{w\}\) falsifies its forward implication at \(w\).
If \(w\) lacks its own loop, take \(V(p)=W\setminus\{w\}\):
all successors satisfy \(p\), but \(w\) does not, falsifying
the reverse implication. Thus validity forces exactly identity.
The box of falsehood is true everywhere exactly when there is no edge.
Finally a loop witnesses \(p\to\Diamond p\) for every valuation;
if \(w\) has no loop, \(V(p)=\{w\}\) falsifies this implication.
These arguments prove both directions of every displayed correspondent,
not merely soundness on the proposed frame classes.
\end{proof}
\end{annexitem}

\begin{annexitem}{annex:p08.src025501.frame-transform}{thm:p08-bvb501-frame-transformations}{Pulling back and extending arbitrary valuations}
\begin{proof}
On a disjoint union, restrict any valuation to each component. Model
truth preservation of component inclusions gives validity at every
union point if each component validates the formula. For a generated
subframe, extend any of its valuations arbitrarily to the original
frame. Generated-submodel truth preservation transfers the original
validity back to every retained point. For a surjective bounded
morphism \(f:F\to F'\), pull any valuation \(V'\) back by
\(V(p)=f^{-1}(V'(p))\). This makes \(f\) a model bounded
morphism. Every \(w'\) has a preimage \(w\); model truth
preservation transfers validity at \(w\) to \(w'\). As the
target valuation was arbitrary, frame validity follows. These arguments
apply formula by formula to any defining set. The three failed-closure
examples in the body therefore rule out even definition by a set, and
the nonempty domain conditions retain the standing frame convention.
\end{proof}
\end{annexitem}

\begin{annexitem}{annex:p08.src025501.ultrafilter}{thm:p08-bvb501-ultrafilter-reflection}{The ultrafilter truth lemma and the order extension}
\begin{proof}
The operator \(l\) preserves \(W\), finite intersections, and
inclusion. For each ultrafilter \(U\), its \(l\)-preimage
\(\mathcal B_U=\{Y:l(Y)\in U\}\) is a filter, possibly
improper when there is no successor. For any \(X\), if
\(l(X)\in U\), every \(R^{\mathrm{ue}}\)-successor contains
\(X\) by definition. If \(l(X)\notin U\), the family
\(\mathcal B_U\cup\{W\setminus X\}\) has the finite
intersection property. Otherwise a finite intersection of its
\(\mathcal B_U\) members lies inside \(X\); preservation of
intersections and inclusion would force \(l(X)\in U\), a
contradiction. Extend this proper filter to an ultrafilter \(U'\).
It is an extension successor of \(U\) and contains \(W\setminus X\),
so it does not contain \(X\). Thus every extension successor contains
\(X\) iff \(l(X)\in U\).

Induction now proves the model truth lemma. Atoms are defined by
membership; complement and finite intersection use the ultrafilter laws,
and the preceding equivalence proves the box case with
\(X\) the original argument truth set. Diamond follows by duality.
The principal ultrafilter \(\pi_w=\{X:w\in X\}\) therefore
keeps precisely the original truths at \(w\). If a frame countervaluation
exists, its extended valuation gives a counterexample at \(\pi_w\).
This proves validity reflection, with its required direction.

For \((\mathbb N,<)\), \(l(X)=\{n:\{j:j>n\}\subseteq X\}\)
is nonempty iff \(X\) is cofinite, and then is cofinite. Hence a
principal ultrafilter at \(n\) accesses exactly the principal
ultrafilters at \(j>n\) and all nonprincipal ultrafilters. A
nonprincipal \(U\) has \(l(X)\in U\) iff \(X\) is cofinite,
so it accesses exactly all nonprincipal ultrafilters, including itself.
The cofinite filter extends to a nonprincipal ultrafilter: containing a
finite set would contradict membership of its cofinite complement.
There are at least continuum many such extensions. Code the nodes of
the countable full binary tree by integers; the node sets on distinct
infinite branches are infinite and have finite intersection. Extend
the cofinite filter with each branch set. Two resulting ultrafilters
cannot coincide, since then they would contain that finite intersection.
This verifies the uncountable cluster in the source's picture.
Every extension point accesses a reflexive cluster point; no point
of \((\mathbb N,<)\) has a reflexive successor. Reflection therefore
rules out modal definability of the asserted class.
\end{proof}
\end{annexitem}

\begin{annexitem}{annex:p08.src025501.frame-so}{thm:p08-bvb501-frame-second-order}{Quantifying over points and all valuation sets}
\begin{proof}
Standard-translation truth preservation applies to each fixed point
and valuation. A valuation of the finitely many occurring atoms is
exactly a tuple of subsets interpreting \(P_1,\ldots,P_n\).
Universal quantification over these tuples and then all points thus
expresses exactly frame validity. There is no restriction to definable
subsets. For \(p\to\Diamond p\), reflexivity makes its
translation true for every subset. Conversely assign \(P=\{w\}\)
and \(x=w\). Its true antecedent demands an accessible point
belonging to that singleton, forcing \(Rww\). As \(w\) was
arbitrary, the universal second-order condition is precisely reflexivity.
\end{proof}
\end{annexitem}

\begin{annexitem}{annex:p08.src025501.loeb-frame}{thm:p08-bvb501-loeb-correspondence}{Transitivity, absence of infinite chains, and compactness}
\begin{proof}
If transitivity fails through \(wRu,uRv,\neg wRv\), make
\(p\) false exactly at \(u,v\). Among \(w\)'s successors
only \(u\) can be false, and \(u\) has the false successor
\(v\). Thus every \(w\)-successor satisfies \(\Box p\to p\),
but \(\Box p\) fails at \(w\), falsifying L\"ob. If there
is an infinite chain \(w_0Rw_1R\cdots\), make \(p\) false
at all \(w_i\), \(i\geq1\), and true elsewhere. Every false
successor of \(w_0\) is a chain point with a further false
successor, while true successors satisfy the implication automatically.
Again the antecedent holds and the conclusion fails at \(w_0\).
Repeated chain points cause no exception to this valuation argument.

Conversely assume transitivity and no infinite chain. If the antecedent
holds at \(w\) but its conclusion fails, some \(w\)-successor
is false for \(p\). Every false successor must have a further false
successor, since otherwise its \(\Box p\to p\) would fail.
Transitivity keeps every such further point a successor of \(w\).
Dependent choice therefore builds an infinite chain, contradiction.
This proves both frame directions in the standing ZFC setting.

If this frame class were elementary, let \(T\) be a first-order
theory defining it. Add constants \(c_i\), \(i<\omega\), and
sentences \(Rc_ic_{i+1}\). Each finite subset is satisfied in a
sufficiently long finite strict linear order, since such orders satisfy
all of \(T\). Compactness gives a model of \(T\) with an
infinite forward chain, contradiction. Thus no first-order theory
defines the class; the finite modal formula's second-order closure
cannot be replaced by an elementary frame definition.
\end{proof}
\end{annexitem}

\begin{annexitem}{annex:p08.src025501.minimum}{thm:p08-bvb501-minimal-fixed-point}{Intersective closure and the correctly restricted return paths}
\begin{proof}
Let \(A\) be the intersection of all pre-fixed points of a monotone
\(G\); the family is nonempty because \(W\) is pre-fixed.
For each such \(P\), monotonicity gives
\(G(A)\subseteq G(P)\subseteq P\), so \(G(A)\subseteq A\).
Then \(G(G(A))\subseteq G(A)\), making \(G(A)\) itself
pre-fixed, and minimality gives \(A\subseteq G(A)\).
Thus \(G(A)=A\), and it is the least fixed point. A positive
first-order closure matrix defines such a monotone operator. A negative
condition is downward closed, so if any predicate satisfies it, the
empty predicate is its least satisfying one. An unsatisfiable negative
condition has no minimum, explaining the existence qualification.

Fix \(xRy\). Let \(R_xuz\) mean \(Rxu\land Ruz\), and
let \(P\) be all points reachable from \(y\) by finitely many
\(R_x\)-steps. It contains the diamond witness \(y\) and is
closed under edges departing from a \(P\)-point that is an
\(x\)-successor. Hence \(x\) satisfies both antecedent conjuncts.
Frame validity forces \(x\in P\), that is, \(R_x^*yx\).
Conversely, if every successor has such a return path, any true
antecedent has a witness \(y\in P\). Its box conjunct propagates
\(P\)-membership along every departure on the prescribed path,
since each departure is an \(x\)-successor. The endpoint \(x\)
therefore belongs to \(P\), proving the axiom. On transitive frames
every ordinary path starting at an \(x\)-successor stays among
\(x\)'s successors until its endpoint, so the restricted return
condition is equivalent to ordinary finite return. On the three-cycle
of the body, \(P=\{y,z\}\) makes the axiom false at \(x\),
despite ordinary return along all three edges. This disproves the
unrestricted simplification without changing the modal axiom.

Finally the two-loop-successor example has antecedent-true predicates
\(\{u\}\) and \(\{v\}\) whose intersection is empty and
fails the diamond. Thus intersection closure applies after choosing
a witness, not to the entire existential antecedent. The least
restricted reachability predicate is the fixed point of
\(G_{x,y}(P)=\{y\}\cup\{z:\exists u(Rxu\land u\in P\land Ruz)\}\).
This monotone first-order operator supplies the claimed least-fixed-
point correspondence construction.
\end{proof}
\end{annexitem}

\begin{annexitem}{annex:p08.src025501.functional}{thm:p08-bvb501-functional-complexity}{One successor gives a linear-size witness}
\begin{proof}
Unravel a satisfying partial-function model directly from its selected
point. Each node has at most one successor, so the tree is a single
path, finite or infinite, even if the original model had a loop.
Its projection is a bounded morphism onto the reachable generated
submodel, retaining truth. Truncate after the input's modal depth,
using the depth-preservation proof above. This gives at most depth
plus one states and retains partial functionality. It does not rely
on quotient filtration preserving the functional frame condition.
A nondeterministic polynomial-time algorithm guesses such a finite
path and the relevant atom valuations and applies polynomial model
checking, proving membership in NP. Propositional SAT reduces by
using the same formula with no modal operators: any satisfying
valuation gives a singleton edge-free partial-function frame, and
any satisfying model supplies a satisfying propositional valuation
at its selected point. Standard propositional NP-hardness proves
NP-completeness. This argument concerns one functional modality;
several independent functional modalities can create branching and
are not covered by the path-size claim.
\end{proof}
\end{annexitem}

\begin{annexitem}{annex:p08.src025501.frame-enumerate}{thm:p08-bvb501-elementary-enumeration}{Fresh valuation predicates reduce validity to first-order consequence}
\begin{proof}
For a modal input \(\psi\), its proposition predicates are fresh
symbols not in the frame sentence \(\alpha\). Frame validity in
the class means that \(\forall x\operatorname{ST}_x(\psi)\)
holds in every expansion of every \(\alpha\)-frame by arbitrary
valuations. This is exactly ordinary first-order consequence
\(\alpha\entails\forall x\operatorname{ST}_x(\psi)\)
in the expanded language. First-order soundness and completeness
identify it with existence of a finite formal proof. Enumerate all
modal formulas and dovetail their first-order proof searches. Every
valid formula is eventually emitted and no invalid one is emitted.
For an effective frame theory enumerate its axioms and proofs using
finitely many of them; completeness gives the same conclusion. This
enumerates positive answers and does not decide negative answers.
An ineffective elementary theory need not admit the axiom/proof
enumeration, so the stated effectiveness boundary is necessary.
\end{proof}
\end{annexitem}

\begin{annexitem}{annex:p08.src025501.universal}{thm:p08-bvb501-universal-expressivity}{Hidden components and total coverage}
\begin{proof}
Take a singleton isolated model with \(p\) false and compare it
at that point with the disjoint union of it and an isolated \(p\)-true
point. The relation pairing just the original point with its copy is
a basic bisimulation, but \(Ep\) is false in the first model and
true in the second. Hence no basic formula can define this operator.
For a total bisimulation, structural induction already handles basic
operators. If \(E\sigma\) is true, choose its witness anywhere
in the first model; total coverage gives a related point in the second,
and induction preserves \(\sigma\). The converse uses coverage
in the other direction. The universal-box case follows by duality.
Finally the stated encodings follow directly from the satisfaction
clause and nonempty domains: a point satisfying \(A\psi\)
exists exactly when a nonempty model satisfies \(\psi\) everywhere.
Conjoining \(\varphi\) specifies the requested local point.
\end{proof}
\end{annexitem}

\begin{annexitem}{annex:p08.src025501.universal-frames}{thm:p08-bvb501-universal-frames}{Pigeonhole valuations and predecessor witnesses}
\begin{proof}
On at most \(n\) points, \(n+1\) nonempty valuation sets have
two intersecting members: select a point from each and apply the finite
pigeonhole principle. Thus the antecedent forces a true disjunct. On
more than \(n\) points, choose \(n+1\) distinct singleton
valuations; all \(Ep_i\) are true and all intersection diamonds
false, so the formula is not frame valid. For predecessors, if every
\(x\) has some \(yRx\), any true \(p\) at \(x\) gives
\(\Diamond p\) at such a \(y\), hence \(E\Diamond p\)
at \(x\). Conversely give \(p\) the singleton valuation
\(\{x\}\). Frame validity forces a \(\Diamond p\) witness
somewhere, which must be a predecessor of \(x\).

Disjoint unions of sufficiently many singleton frames exceed any fixed
positive bound while preserving all basic frame validities. For the
predecessor condition, use \(\mathbb N\) with \(nRm\)
iff \(n>m\). Every number has a predecessor \(n+1\), but
the generated subframe from \(n\) is \(\{0,\ldots,n\}\)
and its point \(n\) has none. The earlier validity-preservation
theorem therefore rules out basic definitions for both classes.
\end{proof}
\end{annexitem}

\begin{annexitem}{annex:p08.src025501.universal-wf}{thm:p08-bvb501-universal-well-founded}{Global induction and finite acyclic countermodels}
\begin{proof}
On a frame with no infinite chain, suppose \(\Box p\to p\)
is true everywhere but \(p\) false at a selected point. Each false
point must have a false successor, since otherwise its true box would
force \(p\). Dependent choice yields an infinite chain, contradiction.
Conversely on an infinite chain \(w_0Rw_1R\cdots\), make
\(p\) false precisely on its point set. Every false point has a
false successor and every true point satisfies \(\Box p\to p\)
automatically. Thus its universal closure holds while \(p\) fails
at \(w_0\). This proves the correspondent without transitivity.

If a basic formula is invalid anywhere, its negation is satisfiable;
the finite rooted tree theorem produces a finite tree satisfying that
negation with immediate-child edges. Such a tree has no infinite
chain. Therefore the formulas valid on all chain-free frames are
exactly those valid on all frames. If basic formulas defined precisely
the chain-free frames, every defining formula would be such an
all-frame validity, and hence their common valid-frame class would
be all frames, contradiction. This covers definition by an arbitrary
set of basic formulas as well as by one.
\end{proof}
\end{annexitem}

\begin{annexitem}{annex:p08.src025501.difference}{thm:p08-bvb501-difference}{A noninjective total bisimulation forgets difference}
\begin{proof}
A global witness is either the current point or a distinct one, which
proves \(E\varphi\leftrightarrow\varphi\lor D\varphi\).
For the converse failure compare a singleton isolated model with two
isolated states, giving all atoms identical values everywhere. Relate
the singleton to both states. This relation is atomically harmonious,
has vacuous basic edge conditions, and covers both domains, so the
universal-invariance theorem makes all formulas using the smaller
language agree at related points. But \(D\topv\) is false in
the singleton and true at either two-point state. No definition of
this difference formula, hence no definition of the difference operator,
is possible in that smaller language. Totality does not entail the
injectivity that would invalidate this counterexample.
\end{proof}
\end{annexitem}

\begin{annexitem}{annex:p08.src025501.hybrid}{thm:p08-bvb501-hybrid-basic}{Two additional inductive cases and singleton witnesses}
\begin{proof}
At a nominal, its singleton semantics agrees exactly with equality to
its named constant; nominal harmony proves invariance. At \(@_i\chi\),
evaluate \(\chi\) at the named point. Translation correctness
follows by capture-free substitution and induction. Named closure
relates the two named points, so induction also proves bisimulation
invariance at this case. All other cases use the earlier modal proof.
Distribution and self-duality hold because \(@_i\) evaluates at one
fixed point, and any valid formula holds there. The equality schemas
express respectively equality of a name to itself, symmetry, transitivity,
and truth replacement at identical named points; each follows directly
from equality of their denotations. If both diamond witnesses satisfy
\(i\), they are the same point, so that point satisfies both
\(p,q\), proving the intersection diamond. With an ordinary atom
\(r\), take a root with two terminal successors, one satisfying
\(r,p\), the other \(r,q\), neither both \(p,q\). This
falsifies the ordinary-atom version and proves the singleton restriction
essential. Different nominals were not assumed to name different points.
\end{proof}
\end{annexitem}

\begin{annexitem}{annex:p08.src025501.hybrid-frames}{thm:p08-bvb501-hybrid-frames}{Quantifying over all choices of names}
\begin{proof}
At the point named \(i\), the four formulas say respectively that
it has no edge to itself, has no two-step return, has no distinct
successor returning to it, and that the two named points are related
in one direction or identical. Elsewhere the first three implications
have false nominal antecedent. A loop falsifies the first by naming
its point; a two-edge cycle falsifies the second by naming its start,
including a repeated-point loop. A cycle between distinct points
falsifies the third by naming one point. Two distinct unrelated points
falsify the fourth by naming them \(i,j\). Conversely the respective
frame properties exclude exactly these witnesses and force validity
for every nominal assignment. This proves both directions without
assuming that every point already has a fixed name in a particular model.
\end{proof}
\end{annexitem}

\begin{annexitem}{annex:p08.src025501.temporal}{thm:p08-bvb501-temporal-translation}{Interval translation and the separating finite diagrams}
\begin{proof}
Induction translates the chosen endpoint and all intermediate argument
truths; the existential and universal quantifiers range over exactly
the points in the satisfaction clauses. Fresh bound variables prevent
capture, so the displayed clauses preserve truth. For the strict-
expressivity claim, use the two transitive closures specified in the body.
The displayed relation has atomic harmony and matching successors:
each left root and the right root has one \(q\)-true and one
\(q\)-false middle successor, and their transitive closures reach
endpoints with the same \(p\)-true pattern. Thus roots match middle
patterns and endpoints,
and every middle state has exactly endpoint successors of the same
atom pattern. Endpoints are terminal. Thus it is a bisimulation.
The left root's selected endpoint has only \(q\)-true intermediate
points; the right root's only \(p\)-endpoint has a \(q\)-false
intermediate point. Hence \(U(p,q)\) differs at related points,
while basic modal invariance forces agreement on every basic formula.
Both frames are finite, transitive and irreflexive, giving the specified
restricted-class nondefinability as well as general nondefinability.
\end{proof}
\end{annexitem}

\begin{annexitem}{annex:p08.src025501.conditional}{thm:p08-bvb501-conditional-clause}{Finite minima and the centered threshold clause}
\begin{proof}
In a finite preorder, every antecedent world \(u\) has a minimal
antecedent world \(v\) at least as preferred as it: repeatedly move
to a strictly better one; finiteness of strict equivalence classes
terminates this process. Every antecedent world at least as preferred
as this minimal \(v\) is in its equivalence class and is also
minimal. If all minima satisfy \(\psi\), this \(v\) is a
threshold witness for \(u\). Conversely for a minimal \(u\),
its threshold witness \(v\) is equivalent to \(u\); hence
\(C_wuv\), and the inner universal clause forces \(\psi(u)\).
This proves the finite equivalence, including the empty-antecedent case.
Consequent weakening preserves every threshold witness and its universal
condition. Antecedent strengthening does not: in a total order with
\(u\) strictly better than \(v\), let \(\varphi\) hold at
both, \(\psi\) only at \(u\), and \(\theta\) only at
\(v\). Then \(\varphi>\psi\) is true while
\((\varphi\land\theta)>\psi\) is false.

With \(uR_w^cv\) defined by \(C_wvu\), the auxiliary diamond
selects just the threshold \(v\), and its box checks exactly
the worlds \(z\) with \(C_wzv\); the universal modality checks
all initial \(u\). Substituting the original argument truth sets
proves the centered modal reading. If the reference changes at a nested
conditional, the relation becomes \(R_u^c\), not \(R_w^c\),
which proves why this is a contextual reading rather than a single-
relation translation of arbitrarily nested formulas.
On \(\mathbb N\) with preference \(C_wuv\) iff \(u\geq v\),
there is no minimal world. Let \(\varphi=\topv\) and
\(\psi=\bottom\). A minima-only test would be vacuously true,
but every threshold's reflexively related \(z=v\) falsifies the
inner clause. This verifies the needed no-limit distinction.
\end{proof}
\end{annexitem}

\begin{annexitem}{annex:p08.src025501.packed}{thm:p08-bvb501-packed-example}{An outer guard retained in the inner matrix}
\begin{proof}
For any witness \(y\) of the outer conjunction, \(Rxy\)
already holds, so conjoining it to the inner antecedent does not
change truth of that universal implication. If \(Rxy\) fails,
neither formula can use that \(y\) as an outer witness. Thus the
two formulas are equivalent under every assignment. In the corrected
inner guard, \(Rxz\), \(Rzy\) and \(Rxy\) cover all three
pairs of \(x,y,z\), and the quantified matrix's free variable
\(z\) is covered. The outer atomic guard covers \(x,y\).
The first inner guard lacks a conjunct jointly containing \(x,y\),
and a conjunction of two atoms is not an atomic guarded-fragment guard.
Binding \(x\) in place of the original outer \(y\) would
change both the free-variable set and its intended evaluation-point
meaning; keeping \(\exists y\) is necessary for this equivalence.
\end{proof}
\end{annexitem}

\begin{annexitem}{annex:p08.src025501.pdl-regular}{thm:p08-bvb501-pdl-regularity}{Singleton valuations, unfolding and induction}
\begin{proof}
For two relations, equality of their diamond truth for every valuation
of \(p\) forces equality of their successor sets: use each singleton
valuation to test each possible endpoint. Applied to the choice and
composition equations, this proves exactly union and sequential
composition; the converse follows from their relational semantics.
Put \(S=R_\pi\), \(T=R_{\pi^*}\). Unfolding gives
\(T=\mathrm{Id}\cup(S\circ T)\), so \(T\) contains identity
and, inductively, every finite \(S\)-path. Thus \(S^*\subseteq T\).
For the reverse inclusion, fix \(x\) and value \(p\) at precisely
the points reachable from \(x\) by finite \(S\)-paths. This set
contains \(x\) and is \(S\)-closed at every point, so
\([T](p\to[S]p)\) holds at \(x\). Segerberg's axiom implies
\([T]p\) there, forcing every \(T\)-successor into that reachable
set. Hence \(T\subseteq S^*\). Conversely genuine finite closure
satisfies unfolding, and induction propagates \(p\) along every
finite path from a \(p\)-true start. Mutual program semantics then
verifies the displayed conditional and loop encodings: tests are
identity steps restricted to their truth sets. In particular adding
the final \(p?\) rejects an execution ending outside \(p\),
which the source's untested repeat expression otherwise permits.

If a first-order formula \(\theta(x,y)\) uniformly defined finite
reachability, add constants \(a,b\), the sentence \(\theta(a,b)\),
and, for each \(n\), the negation of the formula saying there is
an \(R\)-path of length \(n\) from \(a\) to \(b\), including
length zero. Every finite part is satisfied on a chain longer than
all its prohibited lengths, since \(\theta\) must hold at that
chain's endpoints. Compactness yields \(\theta(a,b)\) with no
finite path of any length, contradiction. This is the all-model
undefinability needed for the source's second-order expressivity claim.
\end{proof}
\end{annexitem}

\begin{annexitem}{annex:p08.src025501.mu}{thm:p08-bvb501-mu-constructions}{Monotonicity, fixed points, program translation and invariance}
\begin{proof}
Induction gives simultaneous polarity claims: all-positive free occurrences
make truth monotone in the variable's assigned set, and all-negative
occurrences make it antitone. Boolean complement reverses the claim;
joins, meets and relational predecessor operators preserve it. For
nested fixed-point binders, if parameter values give operators
\(F\leq G\) pointwise, \(\mu G\) is pre-fixed for \(F\),
so \(\mu F\subseteq\mu G\), and \(\nu F\) is post-fixed
for \(G\), so \(\nu F\subseteq\nu G\). These observations
prove the nested-binder induction as well. The earlier minimum theorem
proves the least formula. For the greatest, complement conjugates
\(F\) to monotone \(H(X)=W\setminus F(W\setminus X)\).
Pre-fixed points of \(H\) correspond, with inclusion reversed,
to post-fixed points of \(F\); hence \(W\setminus\mu H\)
is their union and is the greatest fixed point. This proves both the
union semantics and the syntactic dual, with capture-free substitution.

For \(F(X)=[\![\chi]\!]\cup\operatorname{pre}_a(X)\),
the finite-reachability set is fixed, and every pre-fixed set contains
its endpoints and, by induction on path length, all their finite
predecessors. Thus it is the least fixed point. The set where
\(\chi\) holds at every finitely reachable point is fixed under
\(X\mapsto[\![\chi]\!]\cap\operatorname{box}_a(X)\).
Every post-fixed set satisfies \(\chi\) and remains inside itself
under successors, so all its finite paths satisfy \(\chi\).
It is therefore the greatest fixed point. For
\(X\mapsto\operatorname{pre}_a(X)\), every point of the greatest
fixed point has a successor there; dependent choice yields an infinite
run. Conversely the point set on any infinite run is post-fixed, so
its initial point belongs to the greatest fixed point. These arguments
use fresh binder variables absent from \(\chi\).

For invariance, iterate each least fixed point from empty, using the
operator at successor stages and unions at limits; iterate each greatest
fixed point from the full domain, using intersections at limits.
Monotonicity makes these chains monotone, and a set-sized strictly
changing chain cannot continue past an ordinal larger than the number
of subsets of the domain. Stabilization gives the least/greatest fixed
point, since induction bounds its stages by every pre-/post-fixed point.
Across a bisimulation, related points agree at the initial stage and
every successor stage by structural induction on the binder body, with
coherent variable valuations. Limits preserve agreement by membership
in unions or intersections. Choose an ordinal beyond stabilization in
both models to obtain agreement of the final truth sets. This supplies
the binder case, including nested binders, of full formula invariance.

Mutual induction on PDL programs and formulas proves that
\(T_\pi(\eta)\) denotes precisely the predecessors under \(R_\pi\)
of the continuation truth set. Choice is union, composition composes
these predecessor operations in execution order, and a test intersects
with its already translated argument set. Star is the least fixed
point of \(X\mapsto[\![\eta]\!]\cup\operatorname{pre}_\pi(X)\),
just proved to be finite reachability. Freshness excludes dependence
of tests or the continuation on the new binder, while program induction
makes its occurrences positive. This proves all translation clauses
and their target well-formedness. Finally the corrected MSO clause
quantifies all candidate predicates \(P\); its inner universal
condition is exactly \(F(P)\subseteq P\). Requiring \(Px\)
for every such candidate expresses membership in their intersection,
which is precisely least-fixed-point truth. The predicate and world
renaming conditions permit the same argument at nested binders.
\end{proof}
\end{annexitem}

\begin{annexitem}{annex:p08.src025501.interactions}{thm:p08-bvb501-interactions}{Comparing two-step endpoint sets}
\begin{proof}
At \(x\), \(K_i[a]p\) says that every endpoint of
\(I\circ R_a\) from \(x\) satisfies \(p\), and
\([a]K_ip\) says this for \(R_a\circ I\). Inclusion of
the second endpoint set in the first therefore proves the first
implication for every valuation. Conversely, value \(p\) at
exactly the first endpoint set: any point in the second but not the
first falsifies that implication. Interchanging the endpoint sets
proves the other correspondence.
For a first counterexample, take equivalence classes \(\{x\}\),
\(\{y,z\}\), the sole action edge \(xR_ay\), and
\(V(p)=\{y\}\). At \(x\) the first antecedent is true
but its consequent fails at epistemic alternative \(z\) after
the action. For the reverse, take classes \(\{x,y\}\),
\(\{z\}\), the sole action \(yR_az\), and empty \(V(p)\).
At \(x\), \([a]K_ip\) is vacuous but \(K_i[a]p\)
fails through \(y,z\). Both actions are partial functions.
In a product, a horizontal then vertical path and a vertical then
horizontal path have exactly the same independently chosen coordinate
endpoint; this proves equality of the two compositions, hence both
mixed validities.
\end{proof}
\end{annexitem}

\begin{annexitem}{annex:p08.src025501.barcan}{thm:p08-bvb501-barcan}{Domain inclusions and equality-based separating instances}
\begin{proof}
For decreasing domains, the Barcan antecedent checks \(\chi\)
at every accessible world for every individual in \(\delta_w\).
Each individual's domain \(\delta_v\) at an accessible world
is a subset, so its universal conclusion follows. For increasing
domains, the converse antecedent checks all of \(\delta_v\)
at every accessible world, including all of \(\delta_w\),
so each current individual satisfies the boxed conclusion. Other free
variable values remain assigned as before.

If \(wRv\) and \(a\in\delta_v\setminus\delta_w\), assign
a free variable \(z\) the value \(a\) and use \(\chi=(x\ne z)\).
Every current individual is different from \(a\) at every accessible
world, but \(a\) falsifies the universal at \(v\). This falsifies
Barcan without requiring any extra predicate. If instead
\(a\in\delta_w\setminus\delta_v\), use
\(\chi=\exists y(y=x)\). At every world, all of its locally
quantified individuals satisfy this existence formula, so the converse
Barcan antecedent is true. The current individual \(a\) fails
that formula at \(v\), falsifying its boxed conclusion.
These assignments and formulas are permitted by the outer-domain
convention and also handle empty local domains. Necessity is not
inferred from accidental truth under one fixed assignment. Both
inclusions give equality along edges, but impose no equality on
unrelated components.
\end{proof}
\end{annexitem}

\begin{annexitem}{annex:p08.src025501.lindstrom}{thm:p08-bvb501-modal-lindstrom}{Finite records and the relativization compactness contradiction}
\begin{proof}
Fix an abstract formula's finite relevant atom and modality signature.
At depth zero there are finitely many atomic truth records. Recursively,
a depth-\(k+1\) record specifies its atomic record and, for each
modality, the set of successor depth-\(k\) records. There are
finitely many records at every fixed depth. Each has a modal
characteristic formula: its atom diagram, a diamond for each present
successor record, and a box of the disjunction of the present records,
for each modality. Induction verifies this description.

If an abstract formula is bisimulation invariant and has depth bound
\(k\), unravel each model, then truncate its tree after \(k\)
levels. Bisimulation preserves the first step and the depth property
the second. Trees of height at most \(k\) with equal root depth-
\(k\) records are bisimilar: relate equal remaining-depth records
at equal levels; each specified successor record has a mate and
leaves have no successors. Thus the formula depends only on that
finite root record. Disjoin the characteristic formulas of the records
where it is true, using falsehood for an empty set. Finite occurrence
and expansion invariance lift this equivalence back to the ambient
signature. Conversely basic formulas have their depth bound by a
distance/subformula induction and are bisimulation invariant, so
their expressive power satisfies the stated properties.

For the compactness step suppose some abstract \(\chi\) has no
depth bound. If it had a depth bound on its finite relevant signature,
the preceding argument would give a modal equivalent on that signature.
Expansion invariance would extend the equivalence to all larger
signatures, and the modal equivalent would give a depth bound there
as well. Hence arbitrarily deep disagreement already occurs on the
finite relevant signature. Arbitrarily deep cutoffs disagree with their originals;
one orientation occurs for arbitrarily large depths. Replace \(\chi\)
by its negation if necessary so that the cutoff satisfies it and the
original does not. Take fresh \(p\) and the set
\(\{\neg\chi,\operatorname{Rel}(\chi,p)\}\) together with
every finite word of boxes applied to \(p\), including length zero.
Any finite part is satisfiable: select a sufficiently deep disagreeing
model and let \(p\) denote its cutoff domain. Expansion invariance
retains \(\neg\chi\), relativization yields the other conjunct,
and the finite box demands stay inside the cutoff. Compactness supplies
a model of the whole set. Every point finitely reachable from its
selected point satisfies \(p\). The root-generated submodel of
the full model and of its \(p\)-restriction is therefore identical.
Generated inclusion is a bisimulation, so invariance gives
\(\neg\chi\) there from the full model, but \(\chi\)
there from relativization, contradiction. Hence every formula has a
depth bound and the preceding finite-record argument applies.
Basic logic itself has compactness, proved earlier; relativization is
constructed by restricting diamonds to \(p\)-successors, boxes
to \(p\)-successors by implication, and conjoining \(p\) at
the selected point. Structural induction verifies it. This also
establishes the converse direction under the stated closures.
\end{proof}
\end{annexitem}

\begin{annexitem}{annex:p08.src025501.algebra}{thm:p08-bvb501-algebra-completeness}{Normal Boolean operators and the Lindenbaum counteralgebra}
\begin{proof}
Finite join preservation implies monotonicity of \(m\); its Boolean
dual box preserves finite meets and \(1\). Thus Boolean tautologies,
K's implication axiom and theorem necessitation are sound in every
BAO, as is modus ponens. For a normal extension restrict to algebras
validating its theorems. Closure under uniform substitution permits
substituting arbitrary algebraic term values, so its axioms remain
sound under each valuation; the same rule argument gives soundness.

For completeness define \(\varphi\sim_L\psi\) iff
\(L\vdash\varphi\leftrightarrow\psi\). Boolean reasoning
makes this an equivalence congruence for Boolean operations. Theorem
necessitation and K transfer each implication through box, hence through
its dual diamond, so it is a congruence there too. Operations on
classes are therefore well defined. Propositional identities give a
Boolean algebra, and the normal diamond identities give an additive
zero-preserving operator. Under the valuation \(p\mapsto[p]\),
every formula evaluates to its class by induction. Its class is \(1\)
iff it is an \(L\)-theorem. Moreover any algebraic valuation into
this algebra selects formula classes for the finitely many atoms in
an axiom or theorem; a simultaneous uniform substitution represents
those choices, proving value \(1\) for every \(L\)-theorem
under every valuation. Thus the algebra belongs to the specified
semantic class and separates each non-theorem. The possibly trivial
algebra covers an inconsistent normal extension; no relational-frame
representation with unrestricted valuations is assumed.
\end{proof}
\end{annexitem}

\begin{annexitem}{annex:p08.src025501.neighborhood}{thm:p08-bvb501-neighborhood-results}{Principal filters, failures of K/RN, and the polynomial witness selection}
\begin{proof}
In the relational embedding, membership of an argument truth set in
\(N_w\) means it contains every relational successor; this is
exactly box truth. Induction handles atoms and Booleans, and diamond
duality handles the remaining modality. A singleton family containing
only the successor set would test equality, not inclusion, and would
not provide this proof. The four-point example's three truth sets
are \(\{t,u\}\), \(\{u,v,w\}\) and \(\{u,v\}\),
so its indicated memberships falsify K. The singleton empty-family
example falsifies necessitation. Equivalent arguments have equal
truth sets in every model, so membership gives their box equivalence;
it does not give a box of the equivalence, which can fail in that
same empty-family example.

For the size bound, first replace diamonds by their box duals, with
only linear growth. In a satisfying model select its satisfying point.
For every pair of box-argument formulas having different original
truth sets, select one point where they differ. With \(r\)
arguments this gives a set \(S\) of at most
\(1+r(r-1)/2\) points. Restrict atom valuations to \(S\).
At each selected \(w\), put in its new neighborhood exactly the
sets \([\![\sigma]\!]_M\cap S\) for box arguments
\(\sigma\) whose boxes are true originally at \(w\).
If two restricted argument sets coincide, their original sets coincide,
by the distinguishing points chosen. Their original boxes therefore
have the same truth at \(w\). Structural induction now proves
truth preservation on \(S\): Boolean cases restrict truth sets,
and a box's argument set by induction is its restricted original set;
its membership holds exactly when its original box was true, since
no distinct original argument set can spuriously coincide with a
positive one. Thus the input is still satisfied.
The selected model has polynomially many states and at most \(r\)
explicit subsets in each neighborhood, giving a polynomial-size encoding.
Guess this encoding and model-check by bottom-up set evaluation and
finite subset comparison, a polynomial verification. This proves NP
membership. Propositional SAT reduces unchanged and is satisfiable
on a singleton with an arbitrary neighborhood whenever its valuation
is satisfying, giving NP-hardness. This is not a claim about the
PSPACE-complete relational-K problem.
\end{proof}
\end{annexitem}

\begin{annexitem}{annex:p08.src025501.topology}{thm:p08-bvb501-topological-invariance}{Open matching, pullback maps and finite minimal neighborhoods}
\begin{proof}
For a box true at \(w\), choose an open argument-true neighborhood
\(O\). The matching clause gives \(O'\) whose every point has
a related mate in \(O\); induction transfers argument truth to
all of \(O'\), proving box truth at \(w'\). The symmetric
clause proves the converse. Atoms and Booleans are immediate and
diamond follows by duality. For open restriction, any neighborhood
of a retained point can be intersected with the restricting open set;
the result is open in both spaces. This makes the retained identity
relation a topo-bisimulation and gives truth preservation.

For an open continuous \(f\), the image of a source neighborhood
is an open target neighborhood with a preimage mate for every point.
The inverse image of a target neighborhood is open and every one
of its points has its image in that target neighborhood. These prove
the two matching clauses for the graph. Pullback valuation gives
atomic harmony. To see why arbitrary image valuations do not suffice,
map a two-point discrete space onto a singleton. The map is open
and continuous, but a source atom true at only one point has a true
image atom, so its false source point fails atomic harmony.

On a finite space every point has a smallest open neighborhood,
the intersection of all its neighborhoods. Put \(wRu\) iff
\(u\) belongs to that smallest neighborhood. This is a reflexive
transitive relation: every open set containing \(w\) contains
\(u\), and if every open set containing \(u\) contains
\(v\), it also contains \(v\) whenever it contains \(w\).
Topological box truth is exactly relational box truth for this
relation. Modal equivalence of finite models therefore gives a
relational bisimulation by the proved finite-model theorem. Given an
open \(O\) at a related source point, choose the target's smallest
neighborhood; relational back matches each of its points inside the
source's smallest neighborhood, which lies in \(O\). Reverse
matching uses relational forth. Thus it is a topo-bisimulation.
Finally interior preserves finite intersections and the whole domain,
is contained in its argument and is idempotent. These yield K,
necessitation, T and 4 respectively, proving the stated soundness.
\end{proof}
\end{annexitem}

\begin{annexitem}{annex:p08.src025501.topology-def}{thm:p08-bvb501-topological-definability}{All valuations, clopen sets and named singletons}
\begin{proof}
If every subset is open, interior fixes every atom set, proving
\(p\leftrightarrow\Box p\). Conversely let an arbitrary
subset be the valuation of \(p\); validity of its forward
implication puts every point of that subset in its interior. Hence
the subset is open, proving discreteness. For basic undefinability
of connectedness, any allegedly defining set has a countervaluation
on the disconnected two-point discrete space. Open restriction to
the selected singleton preserves that counterexample, yet the
singleton is connected. This contradicts validity of all defining
formulas on all connected spaces.

In any topology \(\operatorname{int}P\subseteq P\subseteq
\operatorname{cl}P\). Thus the universal antecedent
\(A(\Diamond p\to\Box p)\) means
\(\operatorname{cl}P\subseteq\operatorname{int}P\), equivalently
that \(P\) is both open and closed. In a connected space it
must be empty or whole, giving the consequent. A nontrivial clopen
set supplies a valuation making the antecedent true and both global
consequents false. This proves connectedness definability.

Distinct named points make \(@_i\neg j\) true. The two disjuncts
then require an open neighborhood of one excluding the other, exactly
T0. Quantification over all naming choices proves both frame directions.
For T1, \(\Diamond i\) has extension the closure of the named
singleton; the equivalence requires every singleton to equal its
closure, precisely closed singletons. Conversely closed singletons
give every such equivalence. No claim to a local proof of the
source-cited T2 nondefinability or full topological characterization
is made by these elementary correspondent arguments.
\end{proof}
\end{annexitem}
